% Figure: side-by-side schematic of a prokaryotic and a eukaryotic
% cell, with the canonical organelles labelled. Drawn approximately to
% relative scale on a log axis (linear-scale comparison would render
% the prokaryote invisibly small).
\begin{figure}[htbp]
\centering
\begin{tikzpicture}[
  cellwall/.style={very thick, draw=engblack, fill=engblue!8},
  membrane/.style={thick, draw=engred, fill=engblue!12},
  organelle/.style={thick, draw=engamber!90!black, fill=engamber!20},
  nuc/.style={thick, draw=engblue!80!black, fill=engblue!25},
  lbl/.style={font=\scriptsize, anchor=west},
  ttl/.style={font=\small\bfseries, align=center},
]

% Prokaryotic cell (left)
\node[ttl] at (1.6, 4.3) {Prokaryote \\ \scriptsize $\sim 1$--$5\,\mu$m};

% Outer wall (peptidoglycan)
\draw[cellwall] (1.6, 1.6) ellipse (1.5 and 1.0);
% Inner plasma membrane
\draw[membrane] (1.6, 1.6) ellipse (1.35 and 0.85);
% Nucleoid (no envelope)
\draw[densely dashed, very thick, draw=engblue!80!black]
  (1.6, 1.6) ellipse (0.55 and 0.35);
% Ribosomes
\foreach \p in {(0.5,1.4), (0.7,1.9), (2.6,1.8), (2.4,1.4),
                (1.8,1.1), (1.2,2.1), (2.2,2.0)} {
  \fill[engred] \p circle (0.05);
}
% Flagellum
\draw[engblack, thick, decorate, decoration={snake, amplitude=2pt,
  segment length=6pt}] (3.1, 1.5) -- (4.3, 1.5);

% Prokaryote labels (callouts)
\node[lbl] at (-0.8, 0.4) {peptidoglycan wall};
\draw[thin, draw=engblack] (0.5, 0.6) -- (0.7, 1.0);
\node[lbl] at (3.4, 2.6) {nucleoid (no envelope)};
\draw[thin, draw=engblack] (3.4, 2.55) -- (2.0, 1.85);
\node[lbl] at (3.4, 0.7) {ribosomes};
\draw[thin, draw=engblack] (3.4, 0.7) -- (2.4, 1.35);
\node[lbl] at (4.2, 1.5) {flagellum};

% Eukaryotic cell (right)
\node[ttl] at (10.5, 4.3) {Eukaryote \\ \scriptsize $\sim 10$--$100\,\mu$m};
% Plasma membrane (no wall in animal cells)
\draw[membrane] (10.5, 1.6) ellipse (3.0 and 2.2);
% Nucleus with envelope
\draw[nuc] (9.8, 2.2) ellipse (1.0 and 0.7);
% Nucleolus
\fill[engblue!60!black] (9.5, 2.2) ellipse (0.25 and 0.18);
% Mitochondria (3 instances)
\draw[organelle] (11.3, 2.6) ellipse (0.45 and 0.20);
\draw[engamber!90!black, thick] (10.9, 2.55) -- (11.7, 2.65);
\draw[organelle] (8.7, 1.2) ellipse (0.45 and 0.20);
\draw[engamber!90!black, thick] (8.3, 1.15) -- (9.1, 1.25);
\draw[organelle] (11.8, 1.0) ellipse (0.45 and 0.20);
\draw[engamber!90!black, thick] (11.4, 0.95) -- (12.2, 1.05);
% Endoplasmic reticulum (curved tubes)
\draw[thick, draw=engred!80!black] (10.2, 0.9)
  to[out=20, in=200] (11.5, 1.4)
  to[out=20, in=200] (12.2, 1.5);
\draw[thick, draw=engred!80!black] (10.0, 0.6)
  to[out=10, in=180] (11.3, 0.9);
% Golgi (stack of curved discs)
\draw[thick, draw=engamber!80!black] (9.5, 0.5) to[bend right=20] (10.4, 0.5);
\draw[thick, draw=engamber!80!black] (9.5, 0.35) to[bend right=20] (10.4, 0.35);
\draw[thick, draw=engamber!80!black] (9.5, 0.20) to[bend right=20] (10.4, 0.20);
% Cytoskeleton (sketchy filaments)
\draw[thin, draw=enggray] (8.3, 3.0) -- (12.3, 1.7);
\draw[thin, draw=enggray] (8.5, 2.0) -- (12.5, 2.6);
\draw[thin, draw=enggray] (9.5, 3.2) -- (11.8, 0.7);

% Eukaryote labels
\node[lbl, anchor=east] at (8.0, 0.4) {plasma membrane};
\draw[thin, draw=engblack] (8.0, 0.45) -- (8.4, 0.6);
\node[lbl, anchor=east] at (8.0, 2.2) {nucleus};
\draw[thin, draw=engblack] (8.0, 2.2) -- (8.85, 2.2);
\node[lbl, anchor=east] at (8.0, 2.7) {nucleolus};
\draw[thin, draw=engblack] (8.0, 2.65) -- (9.3, 2.2);
\node[lbl, anchor=west] at (12.0, 3.0) {mitochondrion};
\draw[thin, draw=engblack] (12.0, 3.0) -- (11.6, 2.7);
\node[lbl, anchor=west] at (12.6, 1.5) {ER};
\draw[thin, draw=engblack] (12.6, 1.5) -- (12.3, 1.5);
\node[lbl, anchor=west] at (10.6, -0.05) {Golgi};
\draw[thin, draw=engblack] (10.6, -0.05) -- (10.2, 0.2);
\node[lbl, anchor=west] at (12.6, 2.4) {cytoskeleton};
\draw[thin, draw=engblack] (12.6, 2.4) -- (12.3, 2.4);

\end{tikzpicture}
\caption[Prokaryote and eukaryote, side by side]{A prokaryotic cell
(left, drawn approximately $1$--$5\,\si{\micro\meter}$ across) and a
eukaryotic cell (right, drawn approximately $10$--$100\,\si{\micro\meter}$
across). The prokaryote has a peptidoglycan cell wall, a single plasma
membrane, ribosomes free in the cytoplasm, a single circular
chromosome in a nucleoid region without a membrane envelope, and
appendages such as flagella. The eukaryote has a membrane-bounded
nucleus with a nucleolus, mitochondria with their double membrane,
endoplasmic reticulum, Golgi apparatus, and a cytoskeleton that
organises the interior \cite[ch.~1]{text:v6c01-alberts-mbc-2022}.
Scale within each panel is approximately consistent; the two panels
are at incommensurate scales (linear comparison would render the
prokaryote invisible).}
\label{fig:vol06:ch01:prokaryote-eukaryote}
\end{figure}
